%% kapa: exam=1 round=30 year=2019 period=1 form=A subject=real_estate subject_ko=부동산학원론 seq=23 number=103
%% source: Q-Net 감정평가사 기출문제 (cst00309) artlSeq=5205205 file=extracted/1008/5205205/2019년 제30회 감정평가사 1차시험 문제지/2019년 제30회 감정평가사 1차시험 문제지/2019년도 제30회 감정평가사 1차 1교시 A형.hwp
%% conversion: kordoc 4.24.51 -> tools/gen1.py (auto); review=proofread-sonnet; flags=none
\begin{question}{103}
자본환원율에 관한 설명으로 옳지 않은 것은?
\begin{choices}
  \item 자본환원율이란 대상부동산이 장래 산출할 것으로 기대되는 표준적인 순영업소득과 부동산 가격의 비율이다.
  \item 감가상각 전의 순영업소득으로 가치를 추계하는 경우 감가상각률을 제외한 자본환원율을 사용해야 한다.
  \item 할인현금흐름분석법에서는 별도로 자본회수율을 계산하지 않는다.
  \item 부채감당법에 의한 자본환원율은 부채감당률에 저당비율과 저당상수를 곱하여 구한다.
  \item 지분수익률은 매기간 세전현금수지의 현가와 기말지분복귀액의 현가의 합을 지분투자액과 같게 만드는 내부수익률이다.
\end{choices}
\end{question}
\answer{2}
